%! Unit 5: Determinants and Linear Transformations — Summative Assessment — Unit Test (blank)
\documentclass[10pt]{article}
\usepackage{linalg-article}
\usepackage{linalg-boxes}
\newcommand{\xx}{\mathbf{x}}
\newcommand{\ee}{\mathbf{e}}
\newcommand{\T}{^{\mathsf T}}

% Part-divider strip
\newcommand{\parthead}[1]{%
  \vspace{6pt}\noindent
  \begin{headlinebox}{burgundy}{\color{white}\bfseries #1}\end{headlinebox}%
  \vspace{4pt}}

\begin{document}

\pageheader{Unit 5: Determinants and Linear Transformations}{Unit Test}
\namedateperiod

\vspace{0.06in}
\noindent\textbf{Instructions:} Show all work. No notes unless your teacher says so.

% ======================================================================
\parthead{Part A --- Vocabulary (8 pts)}
Write the letter of the best description next to each term.

\vspace{4pt}
\noindent
\begin{minipage}[t]{0.40\textwidth}
\begin{enumerate}[leftmargin=2.2em, itemsep=7pt, label=\arabic*.]
  \item Cofactor expansion \blank{1.2cm}
  \item Singular matrix \blank{1.2cm}
  \item Determinant \blank{1.2cm}
  \item Product rule \blank{1.2cm}
  \item Minor \blank{1.2cm}
  \item Area/volume factor \blank{1.2cm}
  \item Linear transformation \blank{1.2cm}
  \item Orientation \blank{1.2cm}
\end{enumerate}
\end{minipage}%
\hfill
\begin{minipage}[t]{0.57\textwidth}
\small
\begin{description}[leftmargin=1.4em, itemsep=3pt]
  \item[(A)] Computing a determinant along one row: each entry times its signed minor, with the pattern $+\,-\,+\,\dots$
  \item[(B)] A matrix with $\det=0$: its columns are dependent, it collapses space, and it has no inverse.
  \item[(C)] A number computed from a square matrix; its absolute value is the area (or volume) of the box built from the columns.
  \item[(D)] The rule $\det(AB)=\det A\,\det B$: composing transformations multiplies their scaling factors.
  \item[(E)] The smaller determinant left after crossing out one entry's row and column.
  \item[(F)] The quantity $|\det A|$, the number every area (or volume) is multiplied by under the transformation $A$.
  \item[(G)] A map $\xx\mapsto A\xx$ that fixes the origin and sends grid lines to straight, evenly spaced lines.
  \item[(H)] Whether a transformation keeps or flips the handedness of space --- kept when $\det>0$, flipped when $\det<0$.
\end{description}
\end{minipage}

% ======================================================================
\parthead{Part B --- Multiple Choice (12 pts)}
Circle the best answer.

\begin{enumerate}[leftmargin=1.9em, itemsep=9pt]
  \item If $\det A<0$, the transformation $\xx\mapsto A\xx$
    \hfill (A) collapses space \quad (B) reverses orientation (flips) \quad (C) preserves orientation \quad (D) is the identity
  \item Scaling one row of a matrix by $5$ multiplies the determinant by
    \hfill (A) $1$ \quad (B) $5$ \quad (C) $25$ \quad (D) $0$
  \item Compared with $\det A$, the determinant $\det(A\T)$ is
    \hfill (A) equal \quad (B) the opposite sign \quad (C) the reciprocal \quad (D) zero
  \item The columns of a square matrix $A$ are linearly dependent exactly when
    \hfill (A) $\det A=1$ \quad (B) $\det A=0$ \quad (C) $\det A>0$ \quad (D) $\det A<0$
  \item For an invertible matrix $A$, the determinant $\det(A^{-1})$ equals
    \hfill (A) $\det A$ \quad (B) $-\det A$ \quad (C) $1/\det A$ \quad (D) $0$
  \item A shear such as $\begin{bsmallmatrix}1&1\\0&1\end{bsmallmatrix}$ has determinant
    \hfill (A) $0$ \quad (B) $1$ \quad (C) $2$ \quad (D) $-1$
\end{enumerate}

% ======================================================================
\parthead{Part C --- Short Answer \& Computation (35 pts)}
Show your work.

\begin{enumerate}[leftmargin=1.9em, itemsep=12pt]
  \item Compute the determinant of $A=\begin{bsmallmatrix}2&4\\1&1\end{bsmallmatrix}$. Then give the
        \textbf{area} of the parallelogram built from its columns, and say whether $A$
        \textbf{preserves or reverses orientation}.
        \vspace{1.8cm}
  \item Compute the determinant of $B=\begin{bsmallmatrix}2&1&0\\1&2&1\\0&1&2\end{bsmallmatrix}$ by
        \textbf{cofactor expansion} across the top row (show the three $2\times2$ minors).
        \vspace{2.6cm}
  \item Compute the determinant of $C=\begin{bsmallmatrix}1&2&0\\2&5&1\\0&3&5\end{bsmallmatrix}$ by
        \textbf{reducing to triangular form} (elimination), then multiplying the pivots.
        \vspace{2.8cm}
  \item A $3\times3$ matrix $A$ has $\det A=8$. Find each new determinant and name the row rule you used:
        \textbf{(a)} two rows of $A$ are swapped; \textbf{(b)} one row of $A$ is multiplied by $2$;
        \textbf{(c)} $3\times(\text{row }2)$ is added to row $1$; \textbf{(d)} $\det(2A)$.
        \vspace{2.6cm}
  \item Two $3\times3$ matrices have $\det A=6$ and $\det B=2$. Use the product rule to find
        \textbf{(a)} $\det(AB)$, \textbf{(b)} $\det(A^{-1})$, \textbf{(c)} $\det(A\T)$, \textbf{(d)} $\det(A^2)$.
        \vspace{2.4cm}
  \item Photo software applies the transformation $\xx\mapsto A\xx$ with $A=\begin{bsmallmatrix}4&1\\0&2\end{bsmallmatrix}$.
        \textbf{(a)} What is the area of the image of the unit square? \textbf{(b)} A triangle has area $3$;
        what is the area of its image? \textbf{(c)} Does the image keep its orientation?
        \vspace{2.4cm}
  \item Let $M=\begin{bsmallmatrix}3&1\\6&2\end{bsmallmatrix}$. \textbf{(a)} Compute $\det M$.
        \textbf{(b)} Is $M$ invertible? \textbf{(c)} Describe what the transformation $\xx\mapsto M\xx$ does
        to the unit square (what does it collapse to, and why).
        \vspace{2.4cm}
\end{enumerate}

% ======================================================================
\parthead{Part D --- Extended Response (10 pts)}
Explain your reasoning in full sentences.

\begin{enumerate}[leftmargin=1.9em, itemsep=16pt]
  \item Start with $M=\begin{bsmallmatrix}1&2\\2&5\end{bsmallmatrix}$, which has $\det M=1$. Do one
        elimination step, replacing row $2$ by $(\text{row }2)-2(\text{row }1)$, and look at the new matrix.
        \textbf{(a)} Explain why the determinant is still $1$ after this step. \textbf{(b)} Using the picture
        of the box (parallelogram) built from the rows, explain why sliding it sideways this way keeps its
        area the same.
        \vspace{3.0cm}
  \item A transformation $\xx\mapsto A\xx$ sends the unit square to a parallelogram whose sides are the
        columns $A\ee_1$ and $A\ee_2$. \textbf{(a)} Explain why the area of that parallelogram is $|\det A|$,
        and why this makes $|\det A|$ the factor by which \emph{every} area is scaled. \textbf{(b)} Explain
        why $\det A=0$ forces $A$ to have no inverse, using what happens to the unit square.
        \vspace{2.8cm}
\end{enumerate}

\end{document}
